\section*{Заключение}
\addcontentsline{toc}{section}{Заключение}

В ходе выполнения лабораторных работ была реализована программа, которая позволяет строить связный ациклический ориентированный граф в соответствии с распределением Чампернауна, а также выполнять над графом различные алгоритмы: поиск метрик, метод Шимбелла, подсчёт маршрутов, поиск точек сочленения, алгоритм Дейкстры, вычисление максимального потока и потока минимальной стоимости, подсчёт числа остовных деревьев по теореме Кирхгофа, построение минимального остова алгоритмом Борувки, кодирование и декодирование дерева кодом Прюфера с сохранением весов, эвристическую раскраску графа, проверку на эйлеровость и построение эйлерова цикла, а также построение фундаментальной системы разрезов и операции над разрезами. Реализована защита от некорректного пользовательского ввода.

В ходе выполнения лабораторной работы №1 реализовано формирование случайного связного ациклического ориентированного графа с задаваемым пользователем количеством вершин; степени вершин генерируются в соответствии с распределением Чампернауна. Вычислены эксцентриситеты вершин, центр графа и диаметральные вершины. Реализован метод Шимбелла для нахождения минимальных и максимальных путей с ограничением по числу рёбер. Добавлена возможность подсчёта количества маршрутов между двумя заданными вершинами.

В лабораторной работе №2 для сгенерированных графов реализован поиск точек сочленения в неориентированном графе с помощью обхода в глубину. Реализован классический алгоритм Дейкстры для нахождения кратчайшего пути между выбранными вершинами (с проверкой на отрицательные веса).

В лабораторной работе №3 на основе ориентированного графа сгенерированы матрицы пропускных способностей и стоимостей. Найден максимальный поток алгоритмом Эдмондса–Карпа. Вычислен поток минимальной стоимости заданной величины с использованием алгоритма Дейкстры на приведённых стоимостях.

В лабораторной работе №4 для неориентированного графа вычислено число остовных деревьев с помощью матричной теоремы Кирхгофа. Построен минимальный остов алгоритмом Борувки, выполнено его кодирование в код Прюфера с сохранением весов рёбер и последующее декодирование. Реализован улучшенный алгоритм для нахождения минимальной раскраски графа.

В лабораторной работе №5 проверено, является ли граф эйлеровым; при необходимости выполнена модификация графа для достижения чётности степеней всех вершин, после чего построен эйлеров цикл. На основе минимального остовного дерева построена фундаментальная система разрезов, реализована операция симметрической разности для получения произвольных разрезов по выбору пользователя.

\subsection*{Достоинства}
\begin{itemize}
	\item реализация в объектно-ориентированном стиле: удобное разделение функционала (каждое задание в отдельном модуле);
	\item раскраска графа реализована при помощи улучшенного алгоритма последовательного раскрашивания (жадный алгоритм с эвристикой: порядок вершин отсортирован по убыванию их степени);
	\item логирование изменений при модификации графа в эйлеров.
\end{itemize}

\subsection*{Недостатки}
\begin{itemize}
	\item отсутствие возможности задать параметры распределения Чампернауна пользователем;
	\item В программе существуют похожие по структуре функции, например, генерация различных матриц, которые можно оптимизировать и объединить.
	\item отсутствие возможности выбора ориентированности или её отсутствия при генерации графа;
	\item большая часть алгоритмов работает с использованием матрицы смежности графа, и некторые из них можно ускорить (например нахождение фундаментальной системы разрезов) если использовать списки смежности.
\end{itemize}

\subsection*{Масштабирование программы}
\begin{itemize}
	\item добавление новых методов и алгоритмов для работы с графом, например, алгоритм Дейкстры для отрицательных весов;
	\item добавление других распределений для генерации;
	\item добавление возможности задания параметров распределения пользователем.
\end{itemize}